\documentclass[10pt]{article}

\usepackage[T1]{fontenc}
\usepackage{lmodern}
\usepackage[margin=0.76in]{geometry}
\usepackage{microtype}
\usepackage{amsmath,amssymb,bm}
\usepackage{booktabs,tabularx,array}
\usepackage{enumitem}
\usepackage{xcolor}
\usepackage{hyperref}
\usepackage{titlesec}

\hypersetup{
  colorlinks=true,
  linkcolor=blue!55!black,
  citecolor=blue!55!black,
  urlcolor=blue!55!black,
  pdftitle={Radio-Frequency-Interference Rejection for Pulsed NQR},
  pdfauthor={Technical note prepared from the discussion}
}

\setlength{\parindent}{0pt}
\setlength{\parskip}{4pt plus 1pt}
\setlist[itemize]{leftmargin=1.35em,itemsep=1.5pt,topsep=2pt}
\setlist[enumerate]{leftmargin=1.55em,itemsep=1.5pt,topsep=2pt}
\titleformat{\section}{\large\bfseries}{\thesection}{0.55em}{}
\titleformat{\subsection}{\normalsize\bfseries}{\thesubsection}{0.5em}{}
\titlespacing*{\section}{0pt}{8pt}{3pt}
\titlespacing*{\subsection}{0pt}{6pt}{2pt}
\renewcommand{\arraystretch}{1.16}

\newcommand{\B}{\mathbf B}
\newcommand{\cvec}{\mathbf c}
\newcommand{\hvec}{\mathbf h}
\newcommand{\xvec}{\mathbf x}
\newcommand{\E}{\mathbb E}
\newcommand{\NQR}{\mathrm{NQR}}
\newcommand{\RFI}{\mathrm{RFI}}

\title{\vspace{-1.3em}\textbf{Radio-Frequency-Interference Rejection for Pulsed NQR}\\[-1mm]
\large Reference-coil cancellation, vector pickup, and gated acquisition}
\author{}
\date{}

\begin{document}
\maketitle
\vspace{-2.1em}

\begin{abstract}
\noindent
Ambient radio-frequency interference (RFI) is a practical limiting factor for
low-frequency \({}^{14}\mathrm N\) nuclear-quadrupole-resonance (NQR)
experiments.  The difficult case is coherent interference, often from AM
broadcast stations or nearby electronics, that overlaps the NQR detection band.
This note summarizes candidate mitigation methods and proposes a simulation
roadmap for MRSpinDynamics.  The central recommendation is to model RFI as a
vector magnetic field and to support multi-channel reference-coil cancellation
from the beginning.  In pulsed NQR sequences such as SLSE and SORC, the NQR
receiver is only scientifically useful in the gaps between RF pulses, whereas
reference detectors may run continuously; the implementation should therefore
treat cancellation as a masked/gated estimation problem rather than as ordinary
continuous denoising of a single sampled waveform.
\end{abstract}

\section{Why this belongs in the NQR model}

\({}^{14}\mathrm N\) NQR lines for energetic materials and other nitrogenous
solids can lie in the same frequency range as strong environmental emitters.
In the RAND landmine-detection review, the TNT NQR frequencies are reported as
roughly \(790\)--\(900~\mathrm{kHz}\), directly in the AM broadcast band; the
review states that environmental RF interference can overwhelm TNT responses
and that practical NQR systems have used coil design plus a separate RFI
monitoring antenna followed by subtraction \cite{rand2003}.  A later appendix
in the same review calls out commercial AM stations within about
\(20~\mathrm{kHz}\) of the TNT detection frequencies and nearby electromagnetic
noise sources as current NQR limitations \cite{rand2003}.

This is not only a filtering problem.  The desired NQR signal is transient,
weak, sequence-gated, and often detected with matched filters or echo
integration.  A stationary notch filter may reduce a carrier but can also
distort the baseline, remove useful spectral content, or leave time-varying AM
sidebands.  A physically faithful simulator should instead expose the
experimenter to the trade space among passive rejection, reference-channel
subtraction, active cancellation, pulse gating, and detection bias.

\section{Signal model}

Let the primary NQR receive channel after analog front-end conditioning be
\begin{equation}
  y[n] = s_{\NQR}[n]\,m[n]
       + \sum_{q=0}^{Q-1} (g_q * u_q)[n]
       + \eta_y[n],
  \label{eq:scalar-primary}
\end{equation}
where \(s_{\NQR}\) is the desired NQR response, \(m[n]\in\{0,1\}\) is the
receive-valid mask that removes RF-pulse and ringdown intervals, \(u_q\) are
external interference sources, \(g_q\) are their transfer functions into the
primary channel, and \(\eta_y\) is receiver noise.

For a reference-coil system, the measured reference vector is
\begin{equation}
  \xvec[n] =
  \begin{bmatrix}x_1[n]&x_2[n]&\cdots&x_K[n]\end{bmatrix}^{T},
\end{equation}
and the cleaned primary channel is estimated as
\begin{equation}
  \widehat{s}_{\NQR}[n]
    = y[n] - \sum_{k=1}^{K}(h_k*x_k)[n].
  \label{eq:mimo-canceller}
\end{equation}
The filters \(h_k\) may be complex scalars, short FIR filters, frequency-domain
transfer functions, or adaptive filters.  This is the NQR analogue of adaptive
noise cancellation \cite{widrow1975,haykin2002}, RFI reference-antenna
cancellation in radio astronomy \cite{mitchell2010,finger2018}, and
external-sensor EMI subtraction in low-field MRI \cite{srinivas2021}.

\section{The magnetic field is a vector}

A single reference coil measures only one projection of the RFI magnetic field.
For a small loop with effective pickup vector \(\cvec_k\), the idealized
reference voltage is proportional to
\begin{equation}
  x_k(t) \propto
  \frac{d}{dt}\int_{A_k}\B_{\RFI}(\mathbf r,t)\cdot d\mathbf A
  \approx
  \cvec_k^{T}\frac{d\B_{\RFI}(\mathbf r_k,t)}{dt}.
\end{equation}
The primary NQR receive coil has its own pickup vector \(\cvec_y\).  In the
far-field or spatially uniform limit,
\begin{equation}
  y_{\RFI}(t)\approx \cvec_y^{T}\dot{\B}_{\RFI}(t),
  \qquad
  \xvec(t)\approx C^{T}\dot{\B}_{\RFI}(t),
\end{equation}
where \(C=[\cvec_1,\ldots,\cvec_K]\).  If \(C\) spans the component of
\(\dot{\B}_{\RFI}\) that couples into the primary coil, the primary interference
projection can be estimated.  If not, the cancellation is rank-deficient.

For distant AM transmitters, the wavelength is hundreds of meters and the field
may be nearly uniform over the probe.  Three approximately orthogonal
reference coils near the instrument can then be a reasonable first-order model.
For local electronics, shield currents, cable pickup, or re-radiating metal,
the RFI can be near-field and spatially structured.  A single three-axis
reference station may then fail; the next model is multiple vector stations,
with \(K=3P\) channels for \(P\) physical reference locations.  Electric-field
pickup can also enter through cabling and preamplifier imperfections, so the
software model should allow reference channels that are not purely magnetic.

\section{Pulsed acquisition changes the estimation problem}

In SLSE, SORC, and related pulsed NQR sequences, the primary receive channel is
blanked or unusable during transmit pulses and during subsequent ringdown.  The
desired NQR signal is only acquired in gaps.  Reference detectors, by contrast,
can often run continuously because they are physically decoupled from the NQR
transmit/receive coil or can tolerate the pulse feedthrough.

The cancellation problem should therefore be written with masks.  Let
\(\mathcal G\) be the set of valid receive samples in the pulse gaps and
\(\mathcal B\) be a set of samples believed to contain no NQR response, such as
pre-shot baselines, late-time tails, off-resonance shots, or inter-scan
calibration windows.  A fixed FIR canceller can be estimated by
\begin{equation}
  \widehat{\hvec}
    =
  \arg\min_{\hvec}
  \sum_{n\in\mathcal B}
  \left|y[n]-\Phi_x[n]\hvec\right|^2
  + \lambda \|\hvec\|_2^2,
  \label{eq:gated-ls}
\end{equation}
where \(\Phi_x[n]\) is the tapped-delay design row assembled from the
continuous reference channels.  The cleaned NQR data are then computed only on
\(\mathcal G\).  Continuous reference acquisition is still useful because it
allows:
\begin{itemize}
  \item interpolation of the RFI phase and amplitude across transmit blanking;
  \item tracking of AM modulation envelopes and beat notes over the full shot;
  \item estimation from pulse-contaminated reference intervals while excluding
        those intervals from the primary objective;
  \item transfer-function updates between shots without corrupting the NQR
        response in the detection gaps.
\end{itemize}

This masking is essential for simulation.  A method that appears to work on
continuous synthetic FIDs may fail once the primary channel only contributes
short windows separated by high-power RF pulses.

\section{Candidate RFI-rejection methods}

\begin{table}[h!]
\centering
\caption{RFI mitigation methods to represent in MRSpinDynamics.}
\label{tab:methods}
\footnotesize
\begin{tabularx}{\textwidth}{@{}>{\raggedright\arraybackslash}p{0.19\textwidth}>{\raggedright\arraybackslash}X>{\raggedright\arraybackslash}p{0.25\textwidth}@{}}
\toprule
\textbf{Method} & \textbf{Useful model} & \textbf{Main failure mode} \\
\midrule
Passive shielding and cabling hygiene &
Lower the coupled \(g_q\) before digitization; low-field MRI work emphasizes
diagnostic noise measurements, grounding, cabling, and shielding as prerequisites
for reaching coil-noise limits \cite{guallart2025}. &
Does not remove unavoidable broadcast pickup; cannot be studied as a digital
post-processing algorithm alone. \\
Coil balancing or gradiometry &
Common-mode rejection of spatially smooth ambient magnetic fields, at the cost
of changing sample sensitivity. &
Near-field RFI gradients and sample-signal cancellation. \\
Static notch or spectral excision &
Remove known narrowband carriers after acquisition. &
Signal distortion, AM sidebands, ringing, and data loss near the NQR line. \\
Reference-channel least squares or Wiener cancellation &
Estimate Eq.~\eqref{eq:mimo-canceller} from baseline windows, off-resonance shots,
or late-time gaps. &
Reference channels do not span the primary RFI; reference noise is injected. \\
Adaptive LMS/NLMS/RLS cancellation &
Track slowly changing interference amplitude, phase, and transfer functions.
This is the classical adaptive-noise-canceller family \cite{widrow1975,haykin2002}. &
Adaptation can subtract NQR signal if trained on contaminated windows; LMS may
converge too slowly or misadjust under low reference SNR. \\
Multi-reference cancellation &
Use several reference channels or vector stations; radio astronomy analyses show
that two-reference structures avoid some residual-RFI floors that remain in
single-reference cancellers \cite{mitchell2010}. &
More reference noise and more calibration degrees of freedom. \\
AM carrier/envelope tracking &
Represent AM stations as carrier plus sidebands/envelope and track a compact
parametric model. &
Breaks when interference is broadband, nonstationary, or multi-source. \\
Active cancellation coils &
Drive compensation coils to null the field at the receive coil, optionally with
digital post-cleanup. &
Loop stability, actuator saturation, transmit-pulse feedthrough, and accidental
nulling of useful sample signal. \\
\bottomrule
\end{tabularx}
\end{table}

\section{Active cancellation architecture}

Active RFI cancellation adds one or more drive coils with command vector
\(\mathbf a[n]\).  The primary channel becomes
\begin{equation}
  y[n] = s_{\NQR}[n]m[n] + y_{\RFI}[n] + (G_a*\mathbf a)[n] + \eta_y[n].
\end{equation}
The controller chooses \(\mathbf a\) to minimize the primary RFI estimate, using
reference coils and calibration of the actuator-to-sensor transfer functions.
The natural simulation layers are:
\begin{enumerate}
  \item feedforward cancellation, where references predict the required actuator
        waveform without using the primary NQR channel as an error signal during
        receive gaps;
  \item feedback cancellation during baseline windows, with adaptation frozen
        during expected NQR signal windows;
  \item hybrid analog/digital cancellation, where active coils prevent front-end
        saturation and digital reference subtraction removes residuals.
\end{enumerate}
The simulator should explicitly represent actuator limits and closed-loop
stability.  Digital post-processing can repair additive contamination after
linear acquisition, but it cannot recover data lost to LNA compression, ADC
clipping, or transmit-pulse recovery artifacts.

\section{Simulation observables}

The RFI module should report more than a cleaned waveform.  Useful diagnostics
include:
\begin{itemize}
  \item RFI suppression in dB, measured in the primary receive gaps;
  \item matched-filter or echo-integral NQR SNR before and after cancellation;
  \item NQR amplitude and phase bias caused by cancellation;
  \item residual spectral lines and residual autocorrelation;
  \item reference-noise injection into the cleaned primary channel;
  \item conditioning/rank of the reference-channel design matrix;
  \item fraction of samples masked by transmit pulses and ringdown;
  \item clipping or front-end saturation flags before digital cancellation;
  \item false-alarm and detection-probability changes for a chosen detector.
\end{itemize}

\section{Implementation plan for MRSpinDynamics}

\subsection{Phase 1: synthetic RFI and reference channels}

Add a small signal-processing layer that can generate coherent RFI sources and
reference measurements independent of any NQR sequence:
\begin{itemize}
  \item narrowband tones, AM-modulated carriers, chirps, broadband colored noise,
        and impulsive local interference;
  \item vector magnetic-field sources \(\B_{\RFI}(t)\), with optional electric
        pickup channels;
  \item reference coil objects with pickup vectors, locations, bandwidths, noise
        floors, gain/phase errors, saturation limits, and optional NQR leakage;
  \item primary receive pickup vectors coupled to existing NQR orientation and
        receiver conventions.
\end{itemize}

\subsection{Phase 2: masks for pulsed NQR sequences}

Expose receive-valid masks from SLSE, SORC, FID, and full-density-matrix NQR
simulations.  The mask should distinguish transmit pulse, ringdown, valid
receive, and optional baseline-only windows.  Reference channels should be
sampled continuously over the same shot clock.  This allows a cancellation
algorithm to train on \(\mathcal B\), apply on \(\mathcal G\), and ignore
samples where the primary receiver would be blanked.

\subsection{Phase 3: cancellers}

Implement cancellers as pure numerical transforms with explicit fit/apply
separation:
\begin{itemize}
  \item scalar complex subtraction for a single reference channel;
  \item multi-channel ridge least-squares FIR cancellation using Eq.~\eqref{eq:gated-ls};
  \item frequency-domain transfer-function cancellation for stationary RFI;
  \item LMS/NLMS and RLS adaptive cancellers with mask-aware freeze/update
        controls;
  \item parametric AM carrier trackers as a baseline for broadcast-like RFI.
\end{itemize}
Each canceller should return coefficients, cleaned data, predicted RFI,
reference-noise estimates, and mask-aware diagnostics.

\subsection{Phase 4: active cancellation}

Add optional compensation coils after the passive/digital path is validated.
The first implementation can be feedforward with known actuator transfer
functions.  Later versions can add feedback adaptation, stability margins, and
closed-loop interaction between compensation coils, reference sensors, and the
primary receive coil.

\subsection{Phase 5: examples and tests}

Examples should start with intentionally simple cases and then add the
pathologies that matter experimentally:
\begin{enumerate}
  \item a single AM carrier near a \({}^{14}\mathrm N\) line with one reference
        coil;
  \item three-axis reference coils cancelling a time-varying RFI polarization;
  \item rank-deficient one-axis cancellation that leaves a residual in the
        primary coil;
  \item SLSE receive gaps with continuous references and transmit/ringdown masks;
  \item reference-channel NQR leakage showing amplitude bias;
  \item preamplifier or ADC clipping showing when digital cancellation is too
        late;
  \item active feedforward cancellation followed by digital residual cleanup.
\end{enumerate}
Unit tests should verify analytic least-squares recovery in noise-free cases,
known dB suppression in seeded stochastic cases, no update during masked NQR
windows, and preservation of a synthetic NQR matched-filter amplitude when
reference leakage is zero.

\section{Open design questions}

\begin{itemize}
  \item Should RFI live under \texttt{spin\_dynamics.nqr} or in a shared
        \texttt{spin\_dynamics.interference} namespace usable by low-field MRI
        and ESR examples?
  \item Should reference coils use the existing orientation objects, or should
        they have their own lab-frame geometry model?
  \item What is the right default for transmit-pulse feedthrough into reference
        channels: blank, clip, saturate, or preserve for controller tracking?
  \item How should matched-filter detection report uncertainty after adaptive
        cancellation, given injected reference noise and fitted degrees of
        freedom?
  \item How much active-cancellation control theory belongs in the package
        versus being represented as calibrated transfer functions and actuator
        limits?
\end{itemize}

\begin{thebibliography}{9}

\bibitem{rand2003}
J. MacDonald, J. R. Lockwood, J. McFee, T. Altshuler, T. Broach,
L. Carin, R. Harmon, C. Rappaport, W. Scott, and R. Weaver,
\emph{Alternatives for Landmine Detection}, RAND Corporation, MR-1608,
2003.
\url{https://www.rand.org/content/dam/rand/pubs/monograph_reports/MR1608/RAND_MR1608.pdf}.

\bibitem{widrow1975}
B. Widrow, J. R. Glover, J. M. McCool, J. Kaunitz, C. S. Williams,
R. H. Hearn, J. R. Zeidler, E. Dong, and R. C. Goodlin,
``Adaptive noise cancelling: principles and applications,''
\emph{Proceedings of the IEEE}, vol. 63, no. 12, pp. 1692--1716, 1975.

\bibitem{haykin2002}
S. Haykin, \emph{Adaptive Filter Theory}, 4th ed., Prentice Hall, 2002.

\bibitem{mitchell2010}
D. A. Mitchell, J. G. Robertson, and R. J. Sault,
``Alternative adaptive filter structures for improved radio frequency
interference cancellation in radio astronomy,'' arXiv:1004.0938, 2010.
\url{https://arxiv.org/abs/1004.0938}.

\bibitem{finger2018}
R. Finger, F. Curotto, R. Fuentes, R. Duan, L. Bronfman, and D. Li,
``A FPGA-based fast converging digital adaptive filter for real-time RFI
mitigation on ground based radio telescopes,'' arXiv:1805.06376, 2018.
\url{https://arxiv.org/abs/1805.06376}.

\bibitem{srinivas2021}
S. A. Srinivas, S. F. Cauley, J. P. Stockmann, C. R. Sappo,
C. E. Vaughn, L. L. Wald, W. A. Grissom, and C. Z. Cooley,
``External Dynamic InTerference Estimation and Removal (EDITER) for low
field MRI,'' arXiv:2104.08680, 2021.
\url{https://arxiv.org/abs/2104.08680}.

\bibitem{guallart2025}
T. Guallart-Naval and J. Alonso,
``Electromagnetic noise characterization and suppression in low-field MRI
systems,'' arXiv:2507.20413, 2025.
\url{https://arxiv.org/abs/2507.20413}.

\end{thebibliography}

\end{document}
